% !TEX root = ../manuscript.tex
\begin{abstract}
Knowledge-intensive information systems expose process annotations, retrieval checks, entity bindings, graph nodes, and verification records. Curators inspect these artifacts under fixed audit budgets and need reusable records that make the supplied fidelity field explicit while showing dependency, redundancy, bottleneck exposure, and interpretive role. We introduce Structurally-Calibrated Functional Metacognitive Attribution (SC-FMA), where ``metacognitive'' is used operationally for engineering inspection of observable process artifacts. SC-FMA represents artifacts as dependency-aware audit graphs and uses the Structurally-Calibrated Utility (SCU) objective as a representation constraint system. On a locked PRM800K-derived process-annotation distribution (4,417 samples; 34,219 labeled artifacts), the supervised \wstruct{} fidelity field remains strongest (Spearman $\rho=0.611$), while SC-FMA Ridge tracks it with a small loss ($\rho=0.604$) and adds decomposed record fields. A same-supervision Ridge control using graph-derived features alone reaches only $\rho=0.043$; adding position raises the result to $\rho=0.603$, showing that step position rather than graph topology carries most of the recoverable annotation-order signal. Direct graph-necessity checks reach $\rho\approx0$ for TF-IDF topology and $\rho=0.535$ for a mathematical variable-dependency DAG, below a reverse-position baseline of $\rho=0.568$. A post hoc windowed QP diagnostic raises long-trace consistency from $\rho=0.172$ to $0.385$, but remains below the fidelity-tracking reference. Rule-derived retrieval shows that raw fields recover much of the audit-target signal, with SC-FMA adding organization mainly for redundancy and structural over-correction. The evidence supports a bounded audit-record representation layer; human audit usefulness, production knowledge-base validation, and cross-domain transfer remain future work.
\end{abstract}
